\section{Integrated Discussion}

The main contribution of City1 is not merely a gridded map. It is a disciplined framework for producing, validating, and interpreting a calibrated proxy population surface under missing cell-level truth. This framework requires both analytical stages.

Stage 1 shows that a deterministic calibrated proxy population surface can be constructed reproducibly from open-data features, weak supervision, Random Forest learning, and official-total calibration. The baseline validation stack supports structural credibility through transfer testing, within-city robustness, district aggregation evidence, external structural comparison, ablation, and qualitative plausibility. On its own, however, this deterministic surface still risks being over-read, because a single value per cell can make stronger and weaker support conditions look equally authoritative.

Stage 2 addresses that problem directly. The uncertainty-aware extension adds calibrated ensemble intervals, confidence bands, and hotspot-screening classes so that the baseline surface is not interpreted as equally reliable everywhere. Its strongest support comes from hotspot stability and uncertainty-burden separation, with Almaty emerging as the strongest combined case. The mixed interval coverage, mixed or negative external disagreement alignment, and partial district interval evidence do not invalidate the extension, but they constrain it. They make clear that the correct output identity is reliability-aware screening rather than a true census-uncertainty model.

The two stages protect each other conceptually. Stage 1 without Stage 2 risks deterministic over-reading. Stage 2 without Stage 1 would have no stable calibrated baseline to interpret. Together, they form one coherent calibrated and uncertainty-aware proxy framework.

This dependency is the intellectual center of the paper. Stage 1 gives the framework a stable calibrated surface that can be defended as structurally plausible under missing cell-level truth. Stage 2 then prevents that surface from being read as if every cell carried the same support. The paper is strongest when both statements remain true at once.

City-specific interpretation also becomes sharper under this integrated view. Almaty is the strongest combined case because it pairs strong baseline plausibility with the clearest stable hotspot-screening support. Astana remains useful, but it is caution-heavier both in district support and in the uncertainty-aware layer. Shymkent also remains useful but caution-heavy, with weaker district support and a hotspot structure dominated by review-oriented classes. Semey sits in an intermediate position: it retains a small but meaningful stable hotspot subset while also carrying heavier relative uncertainty and a larger caution tail.

This combined interpretation changes how the framework should be used. The output is suitable for exploratory intra-urban comparison, service-planning screening, hotspot-aware triage, and cautious prioritization. It is not suitable for definitive administrative accounting below the support actually available in the evidence stack, and it should not be treated as if every cell were equally validated.

Those use cases are deliberately modest. The framework is useful when the goal is to identify relative concentration, review potentially important areas, or prioritize where more attention should go next. It is not designed to replace census operations, to automate administrative accounting, or to make unsupervised policy decisions.

The integrated framework also extends naturally into a bounded Stage 3 interpretation layer. In that layer, a tool-grounded assistant consumes the calibrated surface, the interval outputs, the confidence bands, the hotspot classes, and the explicit limitation files through controlled read-only tools rather than by generating unsupported claims from free text alone. Its role is evidence-linked explanation, reviewer-safe wording, and cautious runtime interpretation. It is not a new estimator, not a census-recovery system, and not an automated policy engine.
